\documentclass[11pt,a4paper]{article}
% preamble.tex — estilo común de todos los apuntes Froth.
% Cada apunte hace % preamble.tex — estilo común de todos los apuntes Froth.
% Cada apunte hace % preamble.tex — estilo común de todos los apuntes Froth.
% Cada apunte hace \input{preamble} justo después de \documentclass.
\usepackage[margin=2.4cm]{geometry}
\usepackage{xcolor}
\definecolor{litblue}{RGB}{31,79,121}
\definecolor{litgreen}{RGB}{27,94,32}
\definecolor{litorange}{RGB}{230,81,0}
\usepackage{parskip}
\usepackage{titlesec}
\titleformat{\section}{\Large\bfseries\color{litblue}}{\thesection}{0.6em}{}
\titleformat{\subsection}{\large\bfseries\color{litblue}}{\thesubsection}{0.6em}{}
\usepackage{tcolorbox}
\usepackage{fancyhdr}
\pagestyle{fancy}
\fancyhf{}
\fancyhead[L]{\small\color{litblue}Froth — Apuntes de ML}
\fancyhead[R]{\small\thepage}
\renewcommand{\headrulewidth}{0.4pt}
\usepackage[colorlinks=true,linkcolor=litblue,urlcolor=litblue]{hyperref}

% Cabecera de cada apunte: \cabecera{numero}{titulo}
\newcommand{\cabecera}[2]{%
  \begin{tcolorbox}[colback=litblue,coltext=white,colframe=litblue,arc=2mm]
    {\Large\bfseries Apunte #1 — #2}\\[3pt]
    {\small Proyecto Froth · aprender Machine Learning construyendo}
  \end{tcolorbox}\vspace{0.8em}}

% Cajas temáticas
\newtcolorbox{concepto}[1]{colback=blue!4,colframe=litblue,title={Concepto: #1},arc=1.5mm}
\newtcolorbox{practica}{colback=green!5,colframe=litgreen,title={Pruébalo tú (teclado en mano)},arc=1.5mm}
\newtcolorbox{ojo}{colback=orange!8,colframe=litorange,title={Ojo / error típico},arc=1.5mm}
\newtcolorbox{resumen}{colback=gray!8,colframe=gray!60,title={En una frase},arc=1.5mm}

\newcommand{\codigo}[1]{\texttt{#1}}
 justo después de \documentclass.
\usepackage[margin=2.4cm]{geometry}
\usepackage{xcolor}
\definecolor{litblue}{RGB}{31,79,121}
\definecolor{litgreen}{RGB}{27,94,32}
\definecolor{litorange}{RGB}{230,81,0}
\usepackage{parskip}
\usepackage{titlesec}
\titleformat{\section}{\Large\bfseries\color{litblue}}{\thesection}{0.6em}{}
\titleformat{\subsection}{\large\bfseries\color{litblue}}{\thesubsection}{0.6em}{}
\usepackage{tcolorbox}
\usepackage{fancyhdr}
\pagestyle{fancy}
\fancyhf{}
\fancyhead[L]{\small\color{litblue}Froth — Apuntes de ML}
\fancyhead[R]{\small\thepage}
\renewcommand{\headrulewidth}{0.4pt}
\usepackage[colorlinks=true,linkcolor=litblue,urlcolor=litblue]{hyperref}

% Cabecera de cada apunte: \cabecera{numero}{titulo}
\newcommand{\cabecera}[2]{%
  \begin{tcolorbox}[colback=litblue,coltext=white,colframe=litblue,arc=2mm]
    {\Large\bfseries Apunte #1 — #2}\\[3pt]
    {\small Proyecto Froth · aprender Machine Learning construyendo}
  \end{tcolorbox}\vspace{0.8em}}

% Cajas temáticas
\newtcolorbox{concepto}[1]{colback=blue!4,colframe=litblue,title={Concepto: #1},arc=1.5mm}
\newtcolorbox{practica}{colback=green!5,colframe=litgreen,title={Pruébalo tú (teclado en mano)},arc=1.5mm}
\newtcolorbox{ojo}{colback=orange!8,colframe=litorange,title={Ojo / error típico},arc=1.5mm}
\newtcolorbox{resumen}{colback=gray!8,colframe=gray!60,title={En una frase},arc=1.5mm}

\newcommand{\codigo}[1]{\texttt{#1}}
 justo después de \documentclass.
\usepackage[margin=2.4cm]{geometry}
\usepackage{xcolor}
\definecolor{litblue}{RGB}{31,79,121}
\definecolor{litgreen}{RGB}{27,94,32}
\definecolor{litorange}{RGB}{230,81,0}
\usepackage{parskip}
\usepackage{titlesec}
\titleformat{\section}{\Large\bfseries\color{litblue}}{\thesection}{0.6em}{}
\titleformat{\subsection}{\large\bfseries\color{litblue}}{\thesubsection}{0.6em}{}
\usepackage{tcolorbox}
\usepackage{fancyhdr}
\pagestyle{fancy}
\fancyhf{}
\fancyhead[L]{\small\color{litblue}Froth — Apuntes de ML}
\fancyhead[R]{\small\thepage}
\renewcommand{\headrulewidth}{0.4pt}
\usepackage[colorlinks=true,linkcolor=litblue,urlcolor=litblue]{hyperref}

% Cabecera de cada apunte: \cabecera{numero}{titulo}
\newcommand{\cabecera}[2]{%
  \begin{tcolorbox}[colback=litblue,coltext=white,colframe=litblue,arc=2mm]
    {\Large\bfseries Apunte #1 — #2}\\[3pt]
    {\small Proyecto Froth · aprender Machine Learning construyendo}
  \end{tcolorbox}\vspace{0.8em}}

% Cajas temáticas
\newtcolorbox{concepto}[1]{colback=blue!4,colframe=litblue,title={Concepto: #1},arc=1.5mm}
\newtcolorbox{practica}{colback=green!5,colframe=litgreen,title={Pruébalo tú (teclado en mano)},arc=1.5mm}
\newtcolorbox{ojo}{colback=orange!8,colframe=litorange,title={Ojo / error típico},arc=1.5mm}
\newtcolorbox{resumen}{colback=gray!8,colframe=gray!60,title={En una frase},arc=1.5mm}

\newcommand{\codigo}[1]{\texttt{#1}}

\begin{document}
\cabecera{30}{El nacimiento de specter-mineral: tu primer entrenamiento real}

\section{Qué pasó (en 83 segundos)}

Tu RTX 5060 tomó el SPECTER genérico y lo \textbf{afinó} con los $\sim$1,250 pares
título$\leftrightarrow$abstract del corpus maestro: 2 épocas, 160 pasos, lotes de 16.
El resultado vive en \codigo{models/specter-mineral} — un modelo de embeddings que habla
el dialecto del procesamiento mineral. En CPU habrían sido horas; en tu GPU, 83 segundos.

\section{Los conceptos que acabas de usar}

\begin{concepto}{fine-tuning (afinar) $\neq$ entrenar desde cero}
No construimos un modelo nuevo: tomamos los pesos de SPECTER (que ya sabe ``leer papers'')
y los \textbf{empujamos suavemente} hacia tu dominio. Como un geólogo generalista haciendo
una especialización — no vuelve a la primaria.
\end{concepto}

\begin{concepto}{la jugada elegante: auto-supervisión con MNRL}
Entrenar suele exigir etiquetas humanas (``estos dos textos son parecidos''). No etiquetamos
NADA. El truco: \textbf{MultipleNegativesRankingLoss} — en cada lote de 16 papers, el título
de un paper debe quedar más cerca de \emph{su propio} abstract que de los otros 15 del lote.
Positivos gratis (título y abstract del mismo paper hablan de lo mismo) y negativos gratis
(los demás del lote). Un corpus plano se vuelve dataset de entrenamiento sin mover un dedo.
\end{concepto}

\begin{concepto}{el vocabulario del entrenamiento}
\begin{itemize}
  \item \textbf{loss (pérdida)}: el número que mide ``qué tan equivocado va'' el modelo;
        entrenar = empujar los pesos para bajarlo.
  \item \textbf{paso (step)}: procesar un lote y ajustar pesos una vez. Viste ``160 pasos''.
  \item \textbf{época (epoch)}: una pasada completa por todo el dataset. Hicimos 2.
  \item \textbf{warmup}: los primeros pasos con ajustes suavecitos, para no desbaratar lo
        que el modelo ya sabía.
  \item \textbf{mixed precision (amp)}: cuentas en números de menor precisión donde no
        importa — el motivo de parte de la velocidad en GPU.
\end{itemize}
\end{concepto}

\begin{ojo}
\textbf{Los tropiezos también enseñan}: el entrenamiento falló DOS veces antes de arrancar —
faltaban \codigo{datasets} y luego \codigo{accelerate} (el kit del Trainer de Hugging Face).
Lección de ecosistema: en ML moderno ``instalar la librería'' rara vez basta; el stack de
entrenamiento es su propia bestia. Ambas quedaron en requirements.txt para siempre.
Y antes de eso: pip se negó a cambiar torch CPU$\rightarrow$CUDA porque ``la versión ya
estaba satisfecha'' (mismo número, distinta variante) — se resolvió con
\codigo{--force-reinstall --no-deps}.
\end{ojo}

\section{Cómo se juzga si valió la pena (6.6)}

Un modelo nuevo no es mejor por existir. \codigo{froth/compare.py} somete a AMBOS modelos
al mismo examen sobre los 22 topics: mismos corpus, mismo UMAP (semilla fija), granularidad
elegida por DBCV — ¿quién produce clusters más válidos? Además exporta un \textbf{test
ciego}: dos listados de clusters A/B con la clave escondida, para que el experto (tú)
juzgue sin saber quién es quién. Si specter-mineral no gana, eso TAMBIÉN es un resultado
honesto y publicable.

\begin{resumen}
Fine-tuning = empujar pesos existentes hacia tu dominio; MNRL convierte el corpus en
dataset sin etiquetar (tu título debe amar a tu abstract más que a los ajenos); 83s en tu
GPU tras dos tropiezos de dependencias. El veredicto lo dan DBCV en 22 topics y tu test
ciego — no el entusiasmo.
\end{resumen}

\end{document}
